\documentclass[a4paper,11pt]{article}

\usepackage[utf8]{inputenc}
\usepackage[T2A]{fontenc}
\usepackage[english,russian]{babel}
\usepackage[margin=2.5cm]{geometry}
\usepackage{amsmath,amssymb,amsthm,mathtools}
\usepackage{enumitem}
\usepackage{booktabs}
\usepackage{xcolor}
\usepackage[colorlinks=true,linkcolor=blue,urlcolor=blue,citecolor=blue]{hyperref}
\usepackage{url}
\newcommand{\lean}[1]{\mbox{\path{#1}}}

\theoremstyle{plain}
\newtheorem{theorem}{Теорема}[section]
\newtheorem{lemma}[theorem]{Лемма}
\newtheorem{proposition}[theorem]{Утверждение}
\newtheorem{corollary}[theorem]{Следствие}
\theoremstyle{definition}
\newtheorem{definition}[theorem]{Определение}
\theoremstyle{remark}
\newtheorem*{remark}{Замечание}

\newcommand{\R}{\mathbb{R}}
\newcommand{\Z}{\mathbb{Z}}
\newcommand{\Q}{\mathbb{Q}}
\newcommand{\abs}[1]{\left\lvert #1 \right\rvert}
\newcommand{\set}[1]{\left\{ #1 \right\}}
\newcommand{\ip}[2]{\langle #1, #2\rangle}
\newcommand{\Rpp}[1][2]{\R^{#1}_{++}}
\newcommand{\Rp}[1][2]{\R^{#1}_{+}}
\DeclareMathOperator{\conv}{conv}
\DeclareMathOperator{\cone}{cone}
\DeclareMathOperator{\Sp}{Sp}
\DeclareMathOperator{\STAB}{STAB}
\newcommand{\Spm}{\Sp_\star}
\newcommand{\gs}{\mathrel{\succ}}
\newcommand{\e}{\mathbf e}

\title{NP-трудность разрешимости и слабой разрешимости\\ при фиксированных ценах в размерности $3$}
\author{}
\date{11 октября 2026}

\begin{document}
\maketitle

\begin{abstract}
По графу $G$ за полиномиальное время строятся рациональные цены $P\subset\Rpp[3]$ и рациональный вектор выпусков $y$, для которых $y$ неоклассически разрешим при ценах $P$ $\iff$ $y$ слабо разрешим при ценах $P$ $\iff$ $G$ раскрашивается в три цвета. Поэтому проверка разрешимости и проверка слабой разрешимости при $d=3$ NP-трудны относительно полиномиальной сводимости по Карпу. Обе задачи лежат в NP, так что они NP-полны. Все утверждения доказаны в Lean~4 (библиотека \path{NPHardness}, папка \path{np-hardness/lean}).
\end{abstract}

\section{Постановка}

Обозначения взяты из заметок (\path{latex/notes.tex}). Для набора цен $P=\set{p_t}_{t\in[T]}$ множество $S\subset[T]$~--- \emph{достижимый спектр}, если для каждого $t\in S$ есть ненулевой $\xi\in\Rp[3]$ с $\ip{\xi}{p_s-p_t}>0$ при всех $s\notin S$. Если вместо $>0$ стоит $\ge0$, спектр \emph{слабо достижимый}. Семейства таких спектров обозначаются $\Sp(P)$ и $\widehat{\Sp}(P)$. Рассматриваются два \emph{режима} $\star$: строгий и слабый. Запись $u\gs0$ означает $u>0$ в строгом режиме и $u\ge0$ в слабом, $\Spm(P)$~--- это $\Sp(P)$ или $\widehat{\Sp}(P)$.

По теореме о разрешимости из заметок вектор $y\ge0$ неоклассически разрешим при ценах $P\subset\Rpp[3]$ тогда и только тогда, когда $y\in\cone\set{\mathbf 1_S\mid S\in\Sp(P)}$. По критерию слабой разрешимости при $P\ge0$ вектор $y\ge0$ слабо разрешим тогда и только тогда, когда $y\in\cone\set{\mathbf 1_S\mid S\in\widehat{\Sp}(P)}$. В Lean это \lean{neoSolvable_iff_mem_indicatorCone} и \lean{weakSolvable_iff} библиотеки \path{NeoTiling}.

\begin{theorem}[главная]\label{thm:main}
    По графу $G$ без петель с $n$ вершинами и $m$ рёбрами явно строятся не более $100(n+1)(m+1)$ рациональных цен $P\subset\Rpp[3]$ и рациональный вектор $y\ge0$, для которых следующие условия равносильны.
    \begin{enumerate}[label=(\roman*)]
        \item $y$ неоклассически разрешим при ценах $P$;
        \item $y$ слабо разрешим при ценах $P$;
        \item $G$ раскрашивается в три цвета.
    \end{enumerate}
\end{theorem}

\begin{corollary}
    Проверка неоклассической разрешимости и проверка слабой разрешимости при фиксированных рациональных ценах в $\R^3$ NP-полны относительно полиномиальной сводимости по Карпу.
\end{corollary}

\begin{proof}
    Задача о трёхцветной раскраске NP-полна. Построение теоремы~\ref{thm:main} полиномиально: координаты задаются явными формулами размера $O(\log(nm))$ бит. Принадлежность NP: сертификат~--- не более $T$ спектров разложения $y=\sum_km_k\mathbf 1_{S_k}$ (лемма Каратеодори) вместе с коэффициентами. Достижимость каждого $S_k$ проверяется $\abs{S_k}$ задачами линейного программирования, а разложение~--- решением линейной системы.
\end{proof}

Доказательство теоремы~\ref{thm:main} состоит из трёх частей. Гаджет на параболоиде переводит разрешимость в принадлежность многограннику независимых множеств графа $H$, нарисованного на плоскости (раздел~\ref{sec:gadget}). Условие на рисунок проверяется локально по решётке (раздел~\ref{sec:disks}). Граф $H$~--- это нарисованный на решётке аналог $G\square K_3$, а принадлежность точки $s$ многограннику $\STAB(H)$ равносильна трёхцветной раскраске $G$ (разделы~\ref{sec:layout} и~\ref{sec:coloring}).

\section{Гаджет на параболоиде}\label{sec:gadget}

Пусть $\e=(1,1,1)$. Точку $a\in\R^2$ поднимем на параболоид:
\[
    \ell(a):=(a_1,\ a_2,\ -a_1-a_2)+(1+\abs a^2)\,\e.
\]
Для $c\in\R^2$ положим $\xi_c:=(1-4c_1+2c_2,\ 1-4c_2+2c_1,\ 1+2c_1+2c_2)$.

\begin{lemma}[свидетель круга]\label{lem:wit}
    $\ip{\xi_c}{\ell(a)}=3(1+\abs{a-c}^2-\abs c^2)$ и $\ip{\xi_c}{\e}=3$. При $\abs{c_1}+\abs{c_2}\le\frac14$ вектор $\xi_c$ неотрицателен и ненулевой. Поэтому для $q=\ell(m)+\eta\e$
    \[
        \ip{\xi_c}{\ell(a)-q}=3\bigl(\abs{a-c}^2-\abs{m-c}^2-\eta\bigr).
    \]
\end{lemma}

\begin{proof}
    $\ip{\xi_c}{(a_1,a_2,-a_1-a_2)}=-6\ip ac$, а сумма координат $\xi_c$ равна $3$. Неотрицательность: $4\abs{c_1}+2\abs{c_2}\le1$.
\end{proof}

Иначе говоря, $\ip{\xi_c}{\ell(a)-q}>0$ означает, что $a$ лежит вне круга с центром $c$ и квадратом радиуса $\abs{m-c}^2+\eta$.

\begin{lemma}[хорда]\label{lem:chord}
    Пусть $\eta>\frac14\abs{a-b}^2$ и $q=\ell(\frac{a+b}2)+\eta\e$. Тогда нет ненулевого $\xi\ge0$ с $\ip{\xi}{\ell(a)-q}\ge0$ и $\ip{\xi}{\ell(b)-q}\ge0$.
\end{lemma}

\begin{proof}
    $\ell(a)+\ell(b)=2\ell(\frac{a+b}2)+\frac12\abs{a-b}^2\e$. Сумма двух скалярных произведений равна $(\frac12\abs{a-b}^2-2\eta)\ip{\xi}{\e}<0$.
\end{proof}

Пусть $H=(V,E)$~--- граф с рёбрами $\set{u_k,v_k}$, $k\in E$, точки $a_v\in\R^2$ при $v\in V$ и высоты $\eta_k$. \emph{Гаджет} состоит из $1+\abs V+\abs E$ цен:
\[
    p_z:=\tfrac14\e,\qquad p_v:=\ell(a_v),\qquad p_k:=\ell(m_k)+\eta_k\e,\quad m_k:=\tfrac12(a_{u_k}+a_{v_k}).
\]
Это \emph{маркер}, \emph{награды} и \emph{тяжёлые} точки. Для $s\in\R^V$ вектор выпусков: $y_z=1$, $y_k=1$, $y_v=1-s_v$. Многогранник $\STAB(H)$~--- выпуклая оболочка индикаторов независимых множеств.

\begin{definition}[пустые круги]\label{def:valid}
    Гаджет \emph{корректен}, если
    \begin{enumerate}[label=(\alph*)]
        \item точки $a_v$ попарно различны и $\abs{a_{v,1}}+\abs{a_{v,2}}\le\frac18$;
        \item $\eta_k>\frac14\abs{a_{u_k}-a_{v_k}}^2$ при всех $k$;
        \item для каждого ребра $k$ и каждого независимого множества $I$ есть центр $c$ с $\abs{c_1}+\abs{c_2}\le\frac14$, для которого $\abs{a_w-c}^2>\abs{m_k-c}^2+\eta_k$ при всех $w\in I$.
    \end{enumerate}
\end{definition}

Условие (c)~--- главное. Круг может зависеть от $I$ и обязан обходить только награды из $I$. Именно это позволяет рёбрам $H$ пересекаться (раздел~\ref{sec:layout}).

\begin{theorem}[гаджет]\label{thm:gadget}
    Пусть гаджет корректен. Тогда в любом режиме $y\in\cone\set{\mathbf 1_S\mid S\in\Spm(P)}$ тогда и только тогда, когда $s\in\STAB(H)$. При $s\le\mathbf 1$ это означает: $y$ неоклассически разрешим $\iff$ $y$ слабо разрешим $\iff$ $s\in\STAB(H)$.
\end{theorem}

\begin{proof}
    \emph{Маркер.} Координаты $\ell(a)$ не меньше $1-\abs{a_1}-\abs{a_2}\ge\frac78$, а $\eta_k\ge0$ по (b). Поэтому $p_z<p_t$ покоординатно при $t\ne z$. Если $t\in S$ и $z\notin S$, то свидетель $\xi$ индекса $t$ даёт $\ip{\xi}{p_z-p_t}<0$. Значит, каждый непустой слабо достижимый спектр содержит $z$.

    $(\Rightarrow)$ Пусть $y=\sum_Sm_S\mathbf 1_S$, $m_S>0$, $S\in\Spm(P)$. Все $S$ непусты и содержат $z$, поэтому из $y_z=1$ получаем $\sum_Sm_S=1$. Из $y_k=1$ следует, что каждое $S$ содержит все тяжёлые точки. Если обе награды $u_k,v_k$ лежат вне $S$, то свидетель тяжёлой точки $k$ противоречит лемме~\ref{lem:chord}. Значит, награды вне $S$ образуют независимое множество $I_S$, и $s=\mathbf 1-(y_v)_v=\sum_Sm_S\mathbf 1_{I_S}\in\STAB(H)$.

    $(\Leftarrow)$ Пусть $s=\sum_I\mu_I\mathbf 1_I$. Тогда $y=\sum_I\mu_I\mathbf 1_{S_I}$, где $S_I$~--- дополнение наград из $I$. Покажем, что $S_I\in\Sp(P)$, свидетели строгие. Индексы вне $S_I$~--- награды $w\in I$. Для маркера свидетель $\e$: $\ip{\e}{p_w-p_z}>0$. Для награды $v\notin I$ свидетель $\xi_{a_v}$: $\ip{\xi_{a_v}}{p_w-p_v}=3\abs{a_w-a_v}^2>0$ по (a). Для тяжёлой точки $k$ свидетель $\xi_c$ с центром $c$ из (c): $\ip{\xi_c}{p_w-p_k}=3(\abs{a_w-c}^2-\abs{m_k-c}^2-\eta_k)>0$ по лемме~\ref{lem:wit}.

    Последнее утверждение: цены положительны, а при $s\le\mathbf 1$ вектор $y$ неотрицателен. Остаётся применить теорему о разрешимости и критерий слабой разрешимости.
\end{proof}

\section{Пустые круги на решётке}\label{sec:disks}

Пусть вершины $H$ нарисованы в попарно различных точках $\pi_v\in\Z^2$. Для ребра $k$ с концами $A=\pi_{u_k}$, $B=\pi_{v_k}$ и вершины $w$ положим
\[
    L_k:=\abs{B-A}^2,\qquad u_w:=2\pi_w-A-B.
\]
Вектор $u_w$~--- удвоенный вектор из середины ребра в $w$. Вершина $w$ лежит в замкнутом круге с диаметром $AB$, если $\abs{u_w}^2\le L_k$, и на его окружности, если $\abs{u_w}^2=L_k$.

\begin{definition}[решёточный рисунок]\label{def:drawing}
    Рисунок \emph{допустим}, если для каждого ребра $k$
    \begin{enumerate}[label=(\alph*)]
        \item $L_k\le16$;
        \item $\abs{u_w}^2\ge L_k$ при всех $w$: внутри круга с диаметром $AB$ нет вершин;
        \item на окружности, кроме $u_k$ и $v_k$, лежит не больше одной вершины или ровно два конца некоторого ребра.
    \end{enumerate}
\end{definition}

\begin{lemma}[пустые круги]\label{lem:disks}
    Пусть рисунок допустим, $0<\sigma\le\frac1{16}$ и $\sigma(\abs{\pi_{v,1}}+\abs{\pi_{v,2}})\le\frac1{16}$ при всех $v$. Положим $a_v:=\sigma\pi_v$ и $\eta_k:=\sigma^2(\frac14L_k+\frac1{512})$. Тогда гаджет корректен.
\end{lemma}

\begin{proof}
    Условия (a) и (b) определения~\ref{def:valid} очевидны. Проверим (c). Пусть $k$~--- ребро, $I$~--- независимое множество, $C$~--- вершины $I$ на окружности ребра $k$. Множество $I$ содержит не больше одного конца ребра $k$ и не больше одного конца ребра на окружности. Поэтому $\abs C\le2$. Две вершины $w\ne w'$ из $C$ не диаметрально противоположны. Действительно, $u_{w'}=-u_w$ при $w=u_k$ даёт $\pi_{w'}=B$, то есть $w'=v_k$, а оба конца ребра в $I$ не лежат. Если же $w,w'$~--- концы ребра на окружности, они смежны.

    Положим $v:=\sum_{w\in C}u_w\in\Z^2$, $\delta:=\frac1{256}$ и возьмём центр $c:=\sigma(M-\delta v)$, где $M=\frac12(A+B)$. Для $w\in I$ раскроем:
    \[
        \frac{\abs{a_w-c}^2-\abs{\sigma M-c}^2}{\sigma^2}=\frac{\abs{u_w}^2}4+\delta\ip{u_w}{v}.
    \]
    Нужно, чтобы это было больше $\frac14L_k+\frac1{512}$.

    Пусть $w\in C$. Тогда $\abs{u_w}^2=L_k$ и $\ip{u_w}{v}\ge1$. В самом деле, при $C=\set w$ это $L_k\ge1$. При $C=\set{w,w'}$ это $L_k+\ip{u_w}{u_{w'}}$, и это целое число положительно, так как $\abs{u_w+u_{w'}}^2=2(L_k+\ip{u_w}{u_{w'}})>0$. Отсюда добавка не меньше $\delta=\frac1{256}>\frac1{512}$.

    Пусть $w\notin C$. Тогда $\abs{u_w}^2=L_k+j$ с целым $j\ge1$. Вектор $v$~--- сумма не более двух векторов длины $\sqrt{L_k}$, поэтому $\abs v^2\le4L_k\le64$. По неравенству Коши--Буняковского $\delta\abs{\ip{u_w}{v}}\le\delta\sqrt{64(L_k+j)}$. При $L_k\le16$ величина $\frac j4-\delta\sqrt{64(16+j)}$ возрастает по $j$ и при $j=1$ равна $0{,}121\ldots>\frac1{512}$.

    Наконец, $\sigma(\abs{M_1}+\abs{M_2})\le\frac1{16}$ и $\sigma\delta(\abs{v_1}+\abs{v_2})\le\frac1{16}\cdot\frac{16}{256}$, так что $\abs{c_1}+\abs{c_2}\le\frac14$.
\end{proof}

\section{Рисунок графа $G\square K_3$}\label{sec:layout}

Пусть $G$~--- граф без петель на вершинах $0,\dots,n-1$ с рёбрами $\set{e_k,f_k}$, $k=0,\dots,m-1$. Положим $\underline k:=\min(e_k,f_k)$, $\overline k:=\max(e_k,f_k)$ и $N_k:=4(\overline k-\underline k)\ge4$. Цвета $i\in\set{0,1,2}$. Граф $H$ состоит из \emph{строк}, \emph{столбцов}, \emph{треугольников} и \emph{присоединений}.

\begin{itemize}
    \item \textbf{Строка $(v,i)$}: узлы $r_{v,i,j}$, $0\le j\le12m$, в точках $(x_i+2j,\ 8v+2i)$, где $x_1=3$, $x_0=x_2=1$. Рёбра $r_{v,i,j}r_{v,i,j+1}$.
    \item \textbf{Треугольник вершины $v$}: рёбра между первыми узлами $r_{v,0,0}=(1,8v)$, $r_{v,1,0}=(3,8v+2)$, $r_{v,2,0}=(1,8v+4)$.
    \item \textbf{Столбец $(k,i)$}: узлы $c_{k,i,j}$, $0\le j<N_k$, в точках $(X_{k,i},\ 8\underline k+2i+1+2j)$, где $X_{k,i}=24k+8i+8+2[i=1]$. Рёбра $c_{k,i,j}c_{k,i,j+1}$. Столбец заполняет нечётные высоты строго между строками $(\underline k,i)$ и $(\overline k,i)$.
    \item \textbf{Присоединения}: $c_{k,i,0}$ соединён с $r_{\underline k,i,j_*}$, а $c_{k,i,N_k-1}$ с $r_{\overline k,i,j_*}$, где $j_*=12k+4i+4$. Узел $r_{\cdot,i,j_*}$ стоит в точке $x=X_{k,i}+1$.
\end{itemize}

\begin{lemma}[чётность]\label{lem:parity}
    У каждого узла $x+y$ нечётно: у узлов строк $x$ нечётно, $y$ чётно, у узлов столбцов наоборот. Высота $8v+2i$ определяет строку, абсцисса $X_{k,i}$ определяет столбец. Рисунок инъективен.
\end{lemma}

\begin{proof}
    $2i\le4<8$. Числа $X_{k,i}$ дают остатки $8$, $18$, $0$ по модулю $24$ при $i=0,1,2$ и определяют $i$, а затем $k$.
\end{proof}

Строки идут по чётным высотам через нечётные $x$, столбцы~--- по чётным $x$ через нечётные высоты. Поэтому строка пересекает столбец в точке, которая не является узлом. Через неё проходят ребро строки и ребро столбца. Они перпендикулярны, имеют длину $2$ и общую середину. Отдельный гаджет для пересечений не нужен.

\begin{lemma}\label{lem:drawing}
    Рисунок $H$ допустим.
\end{lemma}

\begin{proof}
    (a) Рёбра строк и столбцов имеют $L=4$, присоединения~--- $L=2$, рёбра треугольников~--- $L=8$, $8$ и $16$. Условия (b) и (c) проверим по типам рёбер.

    \emph{Ребро строки или столбца.} Его середина $M$ целая, и у неё $x+y$ чётно. По лемме~\ref{lem:parity} у каждого узла $w$ вектор $\pi_w-M$ ненулевой, так что $\abs{u_w}^2=4\abs{\pi_w-M}^2\ge4$. На окружности лежат четыре соседа точки $M$: два конца ребра и две точки $M\pm\e^\perp$. Для ребра строки это точки над и под серединой, там могут стоять только узлы столбцов. Если стоят оба, они лежат в одном столбце на высотах, отличающихся на $2$, то есть соединены ребром. Для ребра столбца это точки слева и справа от середины, там могут стоять только узлы одной строки, соседние в ней.

    \emph{Присоединение} $A=(X,Y+1)$, $B=(X+1,Y)$ с чётными $X,Y$ (нижнее; верхнее симметрично). У $u_w$ обе координаты нечётны, поэтому $\abs{u_w}^2\ge2=L$. На окружности лежат точки с $u_w\in\set{\pm1}^2$: $A$, $B$, $(X,Y)$ и $(X+1,Y+1)$. У последних двух $x+y$ чётно, узлов там нет.

    \emph{Треугольник} вершины $v$, $h=8v$. Ребро $(1,h)$--$(3,h+2)$: внутри круга лежат точка $(2,h+1)$ (при $x=2$ узлов нет, так как $X_{k,i}\ge8$) и точки с чётным $x+y$. На окружности кроме концов лежат $(3,h)=r_{v,0,1}$ и $(1,h+2)$. Последняя не узел: строка $(v,1)$ начинается с $x=3$. Ребро $(3,h+2)$--$(1,h+4)$ симметрично, на окружности $r_{v,2,1}$. Ребро $(1,h)$--$(1,h+4)$: внутри лежат $(1,h+2)$, точки с $x\in\set{0,2}$ и точки с чётным $x+y$, узлов среди них нет. На окружности кроме концов лежат $(3,h+2)=r_{v,1,0}$ и $(-1,h+2)$, а последняя не узел.
\end{proof}

\section{Лемма о раскраске}\label{sec:coloring}

Узлы $H$ делятся на \emph{копии} и \emph{антикопии}. Каждый узел имеет \emph{хозяина}: вершину $G$ и цвет. Узел $r_{v,i,j}$ принадлежит $(v,i)$, он копия при чётном $j$. Узлы $c_{k,i,0}$ и $c_{k,i,1}$ принадлежат $(\underline k,i)$, это антикопия и копия. Узлы $c_{k,i,j}$ при $j\ge2$ принадлежат $(\overline k,i)$ и являются копиями при чётном $j$. Положим $s_p:=\frac13$ на копиях и $s_p:=\frac23$ на антикопиях.

Все рёбра строк и столбцов, кроме $c_{k,i,1}c_{k,i,2}$, и все присоединения соединяют копию с антикопией одного хозяина. Здесь $j_*$ чётно, а $N_k-1$ нечётно. На этих рёбрах $s_p+s_q=1$. Ребро $c_{k,i,1}c_{k,i,2}$ (\emph{стык}) соединяет копию $(\underline k,i)$ с копией $(\overline k,i)$. Вершины треугольников~--- копии, и на треугольнике $s$ в сумме даёт $1$.

\begin{lemma}[раскраска]\label{lem:coloring}
    $s\in\STAB(H)$ тогда и только тогда, когда $G$ раскрашивается в три цвета.
\end{lemma}

\begin{proof}
    $(\Leftarrow)$ Пусть $\chi$~--- правильная раскраска и $\tau\in\Z/3$. Множество $I_\tau$ состоит из копий хозяев $(v,\chi(v)+\tau)$ и антикопий остальных хозяев. Оно независимо. Ребро копия--антикопия одного хозяина имеет ровно один конец в $I_\tau$. Стык имел бы оба конца в $I_\tau$ только при $\chi(\underline k)=\chi(\overline k)$. В треугольнике в $I_\tau$ лежит ровно одна вершина. Копия лежит в $I_\tau$ при одном значении $\tau$ из трёх, антикопия~--- при двух. Поэтому $\frac13\sum_\tau\mathbf 1_{I_\tau}=s$.

    $(\Rightarrow)$ Пусть $s=\sum_I\mu_I\mathbf 1_I$, $\mu_I>0$. Для клики $K$ (ребро или треугольник) неравенство $\abs{I\cap K}\le1$ верно на всех независимых $I$. Если $\sum_{p\in K}s_p=1$, то оно обращается в равенство на каждом $I$ разложения: среднее чисел $\le1$ с положительными весами равно $1$. Возьмём одно такое $I$. В каждом треугольнике оно содержит ровно один угол $r_{v,\chi(v),0}$, это задаёт $\chi$. На каждом ребре копия--антикопия $I$ содержит ровно один конец. По индукции вдоль строк и столбцов (по связности рёбер копия--антикопия одного хозяина) копия хозяина $(v,i)$ лежит в $I$ тогда и только тогда, когда $i=\chi(v)$. Если $\chi(\underline k)=\chi(\overline k)=i$, то оба конца стыка $c_{k,i,1}c_{k,i,2}$ лежат в $I$, а это противоречит независимости. Значит, $\chi$ правильная.
\end{proof}

\section{Доказательство главной теоремы}

Пусть $H$~--- рисунок раздела~\ref{sec:layout} и $s$~--- точка раздела~\ref{sec:coloring}. Все координаты узлов удовлетворяют $\abs x+\abs y\le24m+8n+8$. Возьмём
\[
    \sigma:=\frac1{16(24m+8n+8)}.
\]
По леммам~\ref{lem:drawing} и~\ref{lem:disks} гаджет с $a_v=\sigma\pi_v$ и $\eta_k=\sigma^2(\frac14L_k+\frac1{512})$ корректен. Цены $P$~--- маркер, награды и тяжёлые точки этого гаджета, $y$~--- его вектор выпусков для $s$. Так как $s\le\mathbf 1$, по теореме~\ref{thm:gadget} (i) и (ii) равносильны условию $s\in\STAB(H)$, а оно по лемме~\ref{lem:coloring} равносильно (iii).

Число цен: $1+\abs V+\abs E$. Узлов в строках $3n(12m+1)$, в столбцах $3\sum_kN_k\le12nm$. Рёбер не больше узлов плюс $6m+3n$. Всего не больше $100(n+1)(m+1)$. Все числа рациональны: $\sigma$ рационально, а $\ell$, середины и $\abs{\cdot}^2$ полиномиальны. \qed

\begin{remark}
    Прежнее доказательство (\path{complexity/complexity.tex}) требовало условия Габриэля: свидетель тяжёлой точки не зависел от $I$. Тогда рёбра $H$ не пересекаются, граф планарен, и для $\STAB$ планарных графов известна только сводимость по Тьюрингу через теоремы Грётшеля--Ловаса--Схрайвера. Свидетели, зависящие от $I$, снимают планарность. В частности, NP-трудность при $d=3$ теперь получается сводимостью по Карпу и без метода эллипсоидов. Результат по Тьюрингу остаётся верным и следует из настоящего.
\end{remark}

\begin{remark}
    Численная проверка (\path{np-hardness/check_layout.py}). Для графов $K_2$, $K_3$, $C_5$, $K_4$, колеса $W_5$ и графа Петерсена (до $5990$ узлов) условия определения~\ref{def:drawing} выполнены на всех рёбрах. Для раскрашиваемых графов три множества $I_\tau$ независимы и в среднем дают $s$. Для $K_4$ и $W_5$ целочисленное программирование подтверждает, что независимого множества со всеми точными равенствами нет.
\end{remark}

\section{Проверка в Lean}

Библиотека \path{NPHardness} (\path{np-hardness/lean}) зависит только от \path{NeoTiling}. Сборка: \texttt{lake build NPHardness}. Аксиомы~--- только стандартные (\lean{propext}, \lean{Classical.choice}, \lean{Quot.sound}).
\begin{center}
\begin{tabular}{@{}lll@{}}
\toprule
Утверждение & Модуль & Lean \\
\midrule
Определение~\ref{def:valid} & \path{Gadget} & \lean{Gadget.Valid} \\
Теорема~\ref{thm:gadget} & \path{Gadget} & \lean{Gadget.mem_cone_iff}, \lean{Gadget.neoSolvable_iff}, \lean{Gadget.weakSolvable_iff} \\
Определение~\ref{def:drawing} & \path{Disks} & \lean{Drawing} \\
Лемма~\ref{lem:disks} & \path{Disks} & \lean{Drawing.gadget_valid} \\
Лемма~\ref{lem:parity} & \path{Layout} & \lean{Graph.pos_shape}, \lean{Graph.pos_injective} \\
Лемма~\ref{lem:drawing} & \path{Layout} & \lean{Graph.short}, \lean{Graph.empty}, \lean{Graph.circle}, \lean{Graph.drawing} \\
Лемма~\ref{lem:coloring} & \path{Coloring} & \lean{Graph.inStab_iff_colorable} \\
Теорема~\ref{thm:main} & \path{Main} & \lean{Graph.neoSolvable_iff_colorable}, \lean{Graph.weakSolvable_iff_colorable}, \\
 & & \lean{Graph.numPrices_le}, \lean{Graph.prices_rat}, \lean{Graph.outputs_rat} \\
\bottomrule
\end{tabular}
\end{center}
Не формализованы NP-полнота задачи о трёхцветной раскраске и полиномиальность построения как алгоритма (в Mathlib нет классов сложности). Формализованы размер экземпляра и рациональность чисел.

\end{document}
